% !TEX root = main.tex
\subsection{Modifying the ZKBoo Protocol}
\label{subsection:overview-lattice-sms-modifying-zkboo}
In the original ZKBoo protocol~\cite{USENIX:GiaMadOrl16},
the computation trace of a circuit $\Circuit$,
treated as a witness $W$,
is decomposed via a randomized function $\decompose(W)$, 
into three views $\{\view_j\}_{j = 1}^{3}$, 
such that any two views reveal no information of $W$,
then committed by the prover as $\{\cmtview_j\}_{j = 1}^{3}$.
The verifier challenges the prover to open two of the three commitments;
consistency of the two opened views certifies correct evaluation.
\ifcameraready
See the full version for more details.
\else
See Appendix~\ref{appendix:zkboo-preliminaries} for details.
\fi

\medskip
\noindent\textbf{Modified decomposition.}
For proving $\RelationLatticeSMS$,
the witness decomposes as $\Wpre$ and $\Wsuf$
(the pre- and post-challenge components of Section~\ref{subsection:transforming-plonk-arithmetization-to-quadratic-constraints}):
$\Wpre$ is determined before the challenges $\alpha, \beta$ are known,
and $\Wsuf$ is determined only after.

Specifically, we run $\decompose(\Wpre)$ to obtain
%views
\emph{pre-views}
$\{\preview_j\}_{j = 1}^{3}$ 
%which we call pre-views
% satisfying $\Wpre = \sum_{j = 1}^{3} \preview_j$,
and treat each $\preview_j$ as the prefix of the full view:
$\view_j = (\preview_j \| \sufview_j)$.
The suffix views $\{\sufview_j\}_{j = 1}^{3}$ are obtained by running
$\decompose(\cdot)$ 
on the combined process of
(i) computing all components of $\Wsuf$ using entries of $\{\preview_j\}_{j = 1}^{3}$ and
(ii) checking system~\eqref{equation:plonk-quadratic-constraints-with-challenges}.
The resulting full views $\{\view_j\}_{j = 1}^{3}$ are thus
a valid ZKBoo decomposition of the extended witness $(\Wpre \| \Wsuf)$.
%\makeblue{Generating the suffix views may additionally }
%use the entries of $\{\preview_j\}_{j = 1}^{3}$ as needed.
%\makeblue{In other words, generating suffix views can be seen as the continuation of the decomposition from decomposing $\Wpre$ after $\alpha, \beta$ are known.}

\medskip
\noindent\textbf{Modified protocol.}
The modified protocol proceeds in three phases.
\begin{enumerate}
\item The prover commits to $\{\preview_j\}_{j = 1}^{3}$,
sending commitments $\{\cmtpreview_j\}_{j = 1}^{3}$ to the verifier.
\item The verifier sends $\alpha, \beta \uniformly \ZZq$
(derived from the random oracle in the non-interactive version).
The prover computes $\{\sufview_j\}_{j = 1}^{3}$
and sends commitments $\{\cmtsufview_j\}_{j = 1}^{3}$.
\item The commitment to $\view_j$ is the pair $(\cmtpreview_j, \cmtsufview_j)$.
The prover and verifier execute the original ZKBoo consistency check
on the requested views.
\end{enumerate}
This structure ensures that $\alpha, \beta$ are unknown to the prover
when it commits to $\witpre$,
preventing the cheating-prover attack identified in
Section~\ref{subsection:proving-sms-function-evaluations}.
See
\ifcameraready the full version
\else Appendix~\ref{subappendix:lattice-sms-modifying-zkboo}
\fi for the full protocol description and security analysis.
